
\documentclass[11pt]{article}

\usepackage[margin=1in]{geometry}
\usepackage[T1]{fontenc}
\usepackage{lmodern}

\usepackage{amsmath, amssymb}

\usepackage{setspace}

\usepackage[most]{tcolorbox}
\usepackage{xcolor}

\tcbset{
  enhanced,
  boxrule=0pt,
  frame hidden,
  breakable
}

\usepackage{fancyhdr}
\pagestyle{fancy}
\fancyhf{}
\setlength{\headheight}{52pt}
\renewcommand{\headrulewidth}{0pt}

\newcommand{\Course}{Math 4610 Spring 2026}
\newcommand{\Instructor}{Homework 3}
\newcommand{\AssignmentType}{Homework} % or Discussion
\newcommand{\AssignmentNumber}{3}

\newcommand{\Contributors}{
Marks, Stephen
}

\fancyhead[L]{\Course\\ \Instructor}
\fancyhead[R]{\Contributors}

\newcommand{\ProblemDivider}[1]{%
  \vspace{1.0em}
  \noindent\textbf{#1}\par
  \vspace{0.35em}
}

\definecolor{problemBack}{RGB}{235,242,250}
\definecolor{problemFrame}{RGB}{70,110,160}
\definecolor{exerciseGray}{gray}{0.96}
\definecolor{exerciseGrayFrame}{gray}{0.45}

\newtcolorbox{problemBox}{
  colback=exerciseGray,
  colframe=exerciseGrayFrame,
  arc=2mm,
  left=8pt,right=8pt,top=8pt,bottom=8pt
}

\newtcolorbox{solutionBox}{
  colback=white,
  colframe=white,
  arc=0mm,
  left=0pt,right=0pt,top=0pt,bottom=0pt
}

\newcommand{\ProblemAnnounce}{\noindent}
\newcommand{\SolutionAnnounce}{\noindent\textbf{Solution.}\par\medskip}

\newcommand{\Exercise}[1]{\textbf{Exercise~#1.}\;}

\newcommand{\R}{\mathbb{R}}
\newcommand{\C}{\mathbb{C}}

\begin{document}

\begin{center}
  {\Large \textbf{\AssignmentType~\AssignmentNumber}}
\end{center}
\vspace{0.75em}


\ProblemDivider{5.B.2}

\begin{problemBox}
\ProblemAnnounce
Suppose $V$ is a complex vector space and $T \in \mathcal{L}(V)$ has no eigenvalues.
Prove that every subspace of $V$ invariant under $T$ is either $\{0\}$ or infinite-dimensional.
\end{problemBox}

\begin{solutionBox}
\SolutionAnnounce
\doublespacing

Suppose $W$ is a $T$-invariant subspace of $V$.

If $W = \{0\}$, then we are done.

Assume $W \neq \{0\}$ and suppose toward a contradiction that $W$ is finite-dimensional.
Because $W$ is finite-dimensional and over $\mathbb{C}$, the operator
$T|_W$ must have an eigenvalue (by Theorem 5.19).

Thus there exists $0 \neq w \in W$ and $\lambda \in \mathbb{C}$ such that
\[
T w = \lambda w.
\]

But then $\lambda$ is an eigenvalue of $T$, contradicting the assumption
that $T$ has no eigenvalues.

Therefore, no nonzero finite-dimensional invariant subspace can exist.
Hence every invariant subspace is either $\{0\}$ or infinite-dimensional.

\end{solutionBox}

\ProblemDivider{5.B.4}

\begin{problemBox}
\ProblemAnnounce
Suppose $\mathbb{F}=\mathbb{C}$, $T\in\mathcal{L}(V)$, $p$ is a nonconstant polynomial,
and $\alpha\in\mathbb{C}$. Prove that $\alpha$ is an eigenvalue of $p(T)$
if and only if $\alpha=p(\lambda)$ for some eigenvalue $\lambda$ of $T$.
\end{problemBox}

\begin{solutionBox}
\SolutionAnnounce
\doublespacing

($\Rightarrow$)
Suppose $\alpha$ is an eigenvalue of $p(T)$.
Then there exists $v \neq 0$ such that
\[
p(T)v = \alpha v.
\]

Let $q(z) = p(z) - \alpha$.
Then $q$ is a nonconstant polynomial and
\[
q(T)v = 0.
\]

Because $q$ factors over $\mathbb{C}$,
\[
q(z) = c(z-\lambda_1)\cdots(z-\lambda_m)
\]
for some $c \neq 0$.

Hence
\[
0 = q(T)v = c(T-\lambda_1 I)\cdots(T-\lambda_m I)v.
\]

Since $v \neq 0$, at least one factor $T-\lambda_k I$
is not injective. Thus $\lambda_k$ is an eigenvalue of $T$.
Because $\lambda_k$ is a zero of $q$,
\[
p(\lambda_k)=\alpha.
\]

($\Leftarrow$)
Suppose $\alpha=p(\lambda)$ for some eigenvalue $\lambda$ of $T$.
Then there exists $v\neq 0$ such that
\[
Tv=\lambda v.
\]

Then
\[
p(T)v=p(\lambda)v=\alpha v.
\]

Hence $\alpha$ is an eigenvalue of $p(T)$.

\end{solutionBox}

\ProblemDivider{5.B.5}

\begin{problemBox}
\ProblemAnnounce
Give an example of an operator on $\mathbb{R}^2$ that shows the result in Exercise 4
does not hold if $\mathbb{C}$ is replaced with $\mathbb{R}$.
\end{problemBox}

\begin{solutionBox}
\SolutionAnnounce
\doublespacing

Let $T \in \mathcal{L}(\mathbb{R}^2)$ be defined by
\[
T(x,y)=(-y,x),
\]
a rotation by $90^\circ$.

Let $p(t)=t^2$.
Then
\[
p(T)=T^2=-I.
\]

Thus $p(T)$ has eigenvalue $-1$.

However, $T$ has no real eigenvalues (it is a rotation in $\mathbb{R}^2$).
Therefore $-1$ is not equal to $p(\lambda)$ for any eigenvalue $\lambda$ of $T$.

\end{solutionBox}

\ProblemDivider{5.B.6}

\begin{problemBox}
\ProblemAnnounce
Suppose $T \in \mathcal{L}(\mathbb{F}^2)$ is defined by
$T(w,z)=(-z,w)$. Find the minimal polynomial of $T$.
\end{problemBox}

\begin{solutionBox}
\SolutionAnnounce
\doublespacing

Observe that
\[
T^2(w,z)=T(-z,w)=(-w,-z)=- (w,z).
\]

Thus
\[
T^2 = -I.
\]

Hence
\[
T^2 + I = 0.
\]

Therefore $T$ satisfies the polynomial
\[
m(z)=z^2+1.
\]

Because $T$ is not a scalar multiple of the identity,
no polynomial of degree $1$ annihilates $T$.
Thus the minimal polynomial is

\[
\boxed{z^2+1}.
\]

\end{solutionBox}

\ProblemDivider{5.C.1}

\begin{problemBox}
\ProblemAnnounce
Prove or give a counterexample:
If $T^2$ has an upper-triangular matrix with respect to some basis,
then $T$ has an upper-triangular matrix with respect to some basis.
\end{problemBox}

\begin{solutionBox}
\SolutionAnnounce
\doublespacing

Counterexample:

Let $T \in \mathcal{L}(\mathbb{R}^2)$ be defined by
\[
T(x,y)=(-y,x).
\]

Then
\[
T^2=-I,
\]
which is upper-triangular with respect to the standard basis.

However, if $T$ had an upper-triangular matrix,
then $T$ would have a real eigenvalue.
But $T$ has no real eigenvalues (it is a rotation).

Therefore the statement is false.

\end{solutionBox}

\ProblemDivider{5.C.2}

\begin{problemBox}
\ProblemAnnounce
Let $A$ and $B$ be upper-triangular matrices with diagonal entries
$\alpha_1,\dots,\alpha_n$ and $\beta_1,\dots,\beta_n$ respectively.
Show:
(a) $A+B$ is upper-triangular with diagonal $\alpha_j+\beta_j$.
(b) $AB$ is upper-triangular with diagonal $\alpha_j\beta_j$.
\end{problemBox}

\begin{solutionBox}
\SolutionAnnounce
\doublespacing

(a) For $j>k$, we have $A_{jk}=B_{jk}=0$.
Thus
\[
(A+B)_{jk}=A_{jk}+B_{jk}=0.
\]

For diagonal entries,
\[
(A+B)_{jj}=A_{jj}+B_{jj}=\alpha_j+\beta_j.
\]

Thus $A+B$ is upper-triangular.

(b) By matrix multiplication,
\[
(AB)_{jj}=\sum_{r=1}^n A_{jr}B_{rj}.
\]

If $r\neq j$, then either $j>r$ (so $A_{jr}=0$)
or $r>j$ (so $B_{rj}=0$).
Hence all such terms vanish and

\[
(AB)_{jj}=A_{jj}B_{jj}=\alpha_j\beta_j.
\]

If $j>k$, then
\[
(AB)_{jk}=\sum_{r=1}^n A_{jr}B_{rk}.
\]

For $r\le j-1$, $A_{jr}=0$.
For $r\ge j$, since $j>k$, we have $r>k$,
so $B_{rk}=0$.
Hence all terms vanish and $(AB)_{jk}=0$.

Thus $AB$ is upper-triangular with diagonal
$\alpha_1\beta_1,\dots,\alpha_n\beta_n$.

\end{solutionBox}

\end{document}

